\section{Résultats : une efficacité dépendant du réglage et du critère}
\label{sec:results}

Les valeurs proviennent exclusivement des CSV de l'exécution étendue
archivée. Les tableaux ciblés sont exportés automatiquement depuis
\texttt{summary.csv} ; l'arrondi est appliqué seulement à l'affichage.
Les figures PDF existantes sont réutilisées sans régénération. Les
observations suivantes portent sur la grille étudiée, et non sur une
supériorité universelle de Perona--Malik.

\subsection{Comparaison au cas fixé}

La figure~\ref{fig:comparison} compare la vérité terrain, l'observation
commune et les trois sorties au même checkpoint. Les plateaux de la
chaleur deviennent plus réguliers, mais les frontières du rectangle sont
élargies. Avec $K=0{,}10$, la PM exponentielle conserve des pixels isolés
et des fluctuations visibles dans les régions constantes. La rationnelle
atténue davantage ces fluctuations tout en conservant des transitions
visuellement plus localisées. C'est une observation de ce cas, pas une
preuve de préservation exacte des contours.

\begin{figure}[htbp]
 \centering
 \includegraphics[width=\linewidth]{comparison.pdf}
 \caption{Cas individuel préfixé : \texttt{geometric\_shapes},
 $\sigma=0{,}10$, graine~$0$, $n=20$, $t=4$ ; $K=0{,}10$ pour les deux PM,
 $\dx=\dy=1$, $\dt=0{,}20$. Échelle d'affichage commune $[0,1]$,
 sans troncature des tableaux évalués. L'exponentielle conserve ici des
 fluctuations que la rationnelle atténue davantage ; aucune bande statistique.}
 \label{fig:comparison}
\end{figure}

Le profil de la figure~\ref{fig:profile} traverse le rectangle à la
ligne~$42$ (indices commençant à zéro). Contrairement aux images en gris,
il affiche les intensités non tronquées : les pics résiduels et la largeur
des transitions sont directement visibles. La figure~\ref{fig:gradients}
complète cette lecture par une visualisation de
$\sqrt{(D_xu)^2+(D_yu)^2}$ calculée avec \texttt{numpy.gradient}.
Ce gradient de visualisation n'est ni une métrique nouvelle ni le gradient
aux interfaces utilisé pour calculer les conductances directionnelles.
L'échelle colorée est commune ; son maximum est le plus grand des
$99$\textsuperscript{e} percentiles des cinq cartes, de sorte que certaines fortes
valeurs peuvent saturer à l'affichage sans être modifiées numériquement.

\begin{figure}[htbp]
 \centering
 \includegraphics[width=0.95\linewidth]{intensity_profile.pdf}
 \caption{Profil horizontal, ligne~$42$, du même cas individuel
 \texttt{geometric\_shapes}, $\sigma=0{,}10$, graine~$0$, $n=20$, $t=4$,
 $K=0{,}10$, $\dx=\dy=1$, $\dt=0{,}20$. Intensités non tronquées :
 élargissement des transitions par la chaleur et pics résiduels de la PM
 exponentielle. Il ne s'agit pas d'une moyenne sur les graines.}
 \label{fig:profile}
\end{figure}
\begin{figure}[htbp]
 \centering
 \includegraphics[width=\linewidth]{gradient_maps.pdf}
 \caption{Cartes de gradient du cas individuel \texttt{geometric\_shapes},
 $\sigma=0{,}10$, graine~$0$, $n=20$, $t=4$, $K=0{,}10$,
 $\dx=\dy=1$, $\dt=0{,}20$. Échelle commune, plafonnée pour l'affichage
 par les percentiles décrits dans le texte. Les réponses hors des frontières
 illustrent des fluctuations résiduelles ; ce n'est pas une mesure
 quantitative de qualité des contours.}
 \label{fig:gradients}
\end{figure}

Le tableau~\ref{tab:fixed-case}, lui, rapporte les dix graines.
La PM exponentielle améliore le PSNR par rapport à la chaleur
($28{,}3449$ contre $23{,}6482$~dB), tout en donnant un SSIM plus faible
($0{,}6254$ contre $0{,}8017$). L'erreur quadratique globale et les
statistiques structurelles locales ne classent donc pas nécessairement
les sorties de la même manière. La PM rationnelle avec ce $K$ atteint
ici les meilleurs scores des trois solveurs sélectionnés, sans que cette
comparaison à $K$ fixé désigne le meilleur réglage de toute la grille.

\begin{table}[htbp]
 \centering
 \caption{Comparaison agrégée fixée : \texttt{geometric\_shapes},
 $\sigma=0{,}10$, $n=20$, $t=4$, $K=0{,}10$ pour PM ; bruitée à $n=0$.
 Moyenne $\pm$ écart-type échantillonnal sur dix graines, non intervalle
 de confiance. Valeurs exportées depuis \texttt{summary.csv}.}
 \label{tab:fixed-case}
 % Generated from archived summary.csv; rounding applies only to display.
\begin{tabular}{lll}
\toprule
Méthode & PSNR (dB) & SSIM \\
\midrule
Bruitée & $19.9989 \pm 0.0458$ & $0.2168 \pm 0.0013$ \\
Chaleur & $23.6482 \pm 0.0432$ & $0.8017 \pm 0.0023$ \\
PM exp. & $28.3449 \pm 0.2351$ & $0.6254 \pm 0.0137$ \\
PM rat. & $34.6530 \pm 0.2974$ & $0.9292 \pm 0.0027$ \\
\bottomrule
\end{tabular}

\end{table}

\subsection{Toutes les valeurs de \texorpdfstring{$K$}{K} à budget commun}

Le tableau~\ref{tab:common-n20} conserve les huit variantes PM à $n=20$.
L'exponentielle avec $K=0{,}025$ reste proche de l'observation brute :
PSNR $20{,}1269$~dB, SSIM $0{,}2196$, contre $19{,}9989$~dB et $0{,}2168$
pour le bruit. Des différences dues au bruit dépassent elles aussi le
seuil, et leur diffusion devient très faible. Un petit $K$ ne distingue
donc pas automatiquement contour utile et fluctuation indésirable.
Avec $K=0{,}05$, l'exponentielle n'atteint encore que $21{,}1646$~dB.
Ces cas défavorables restent inclus, plutôt que supprimés au profit d'une
seule image flatteuse.

\begin{table}[htbp]
 \centering
 \caption{Toutes les méthodes à budget commun pour
 \texttt{geometric\_shapes}, $\sigma=0{,}10$, $n=20$, $t=4$ ;
 ligne bruitée de référence à $n=0$. Moyennes $\pm$ écarts-types
 échantillonnaux sur dix graines. $K=0{,}025$ est conservé exactement
 comme paramètre, sans arrondi en $0{,}03$.}
 \label{tab:common-n20}
 % Generated from archived summary.csv; rounding applies only to display.
\begin{tabular}{llrll}
\toprule
Méthode & $K$ & $n$ & PSNR (dB) & SSIM \\
\midrule
Chaleur & -- & 20 & $23.6482 \pm 0.0432$ & $0.8017 \pm 0.0023$ \\
Bruitée & -- & 0 & $19.9989 \pm 0.0458$ & $0.2168 \pm 0.0013$ \\
PM exp. & 0.025 & 20 & $20.1269 \pm 0.0483$ & $0.2196 \pm 0.0013$ \\
PM exp. & 0.05 & 20 & $21.1646 \pm 0.0617$ & $0.2455 \pm 0.0018$ \\
PM exp. & 0.1 & 20 & $28.3449 \pm 0.2351$ & $0.6254 \pm 0.0137$ \\
PM exp. & 0.2 & 20 & $35.7056 \pm 0.3280$ & $0.9378 \pm 0.0026$ \\
PM rat. & 0.025 & 20 & $23.7556 \pm 0.0893$ & $0.3276 \pm 0.0032$ \\
PM rat. & 0.05 & 20 & $33.3624 \pm 0.2979$ & $0.8348 \pm 0.0090$ \\
PM rat. & 0.1 & 20 & $34.6530 \pm 0.2974$ & $0.9292 \pm 0.0027$ \\
PM rat. & 0.2 & 20 & $26.9975 \pm 0.1527$ & $0.8667 \pm 0.0033$ \\
\bottomrule
\end{tabular}

\end{table}

\subsection{Temps de diffusion et sur-lissage}

Les figures~\ref{fig:curves025}--\ref{fig:curves10} couvrent les six cas.
Le suffixe image/bruit des fichiers identifie les agrégats sur dix
graines. Le fichier \texttt{metric\_curves.pdf} sans suffixe, qui concerne
uniquement la graine représentative, n'est pas utilisé comme courbe
agrégée. Chaque méthode dispose des mêmes checkpoints ; les bandes
expriment un écart-type, non une incertitude sur une population d'images.

\begin{figure}[p]
 \centering
 \includegraphics[width=0.98\linewidth]{metric_curves_geometric_shapes_sigma0p025.pdf}
 \par\medskip
 \includegraphics[width=0.98\linewidth]{metric_curves_ramps_and_edges_sigma0p025.pdf}
 \caption{PSNR et SSIM agrégés pour $\sigma=0{,}025$ :
 \texttt{geometric\_shapes} en haut, \texttt{ramps\_and\_edges} en bas.
 Moyennes sur dix graines ; bandes $\pm$ un écart-type échantillonnal.
 Chaleur et deux conductances PM, $K\in\{0{,}025;0{,}05;0{,}10;0{,}20\}$,
 $\dx=\dy=1$, $\dt=0{,}20$, $n=0,1,2,5,10,20,40,80$.
 Au faible bruit, un lissage plus intense peut dégrader la fidélité
 quadratique alors que les fluctuations locales diminuent.}
 \label{fig:curves025}
\end{figure}
\begin{figure}[p]
 \centering
 \includegraphics[width=0.98\linewidth]{metric_curves_geometric_shapes_sigma0p05.pdf}
 \par\medskip
 \includegraphics[width=0.98\linewidth]{metric_curves_ramps_and_edges_sigma0p05.pdf}
 \caption{PSNR et SSIM agrégés pour $\sigma=0{,}05$ :
 \texttt{geometric\_shapes} en haut, \texttt{ramps\_and\_edges} en bas.
 Dix graines, moyenne $\pm$ écart-type échantillonnal, non intervalle
 de confiance. Chaleur et PM exponentielle/rationnelle, quatre $K$
 prédéfinis de $0{,}025$ à $0{,}20$ ; $\dx=\dy=1$, $\dt=0{,}20$,
 checkpoints $0,1,2,5,10,20,40,80$. Les évolutions dépendent de la
 conductance et du seuil, et ne sont pas monotones pour toutes les variantes.}
 \label{fig:curves05}
\end{figure}
\begin{figure}[p]
 \centering
 \includegraphics[width=0.98\linewidth]{metric_curves_geometric_shapes_sigma0p1.pdf}
 \par\medskip
 \includegraphics[width=0.98\linewidth]{metric_curves_ramps_and_edges_sigma0p1.pdf}
 \caption{PSNR et SSIM agrégés pour $\sigma=0{,}10$ :
 \texttt{geometric\_shapes} en haut, \texttt{ramps\_and\_edges} en bas.
 Moyennes et bandes $\pm$ un écart-type échantillonnal sur dix graines.
 Chaleur et huit PM avec $K=0{,}025;0{,}05;0{,}10;0{,}20$ ;
 $\dx=\dy=1$, $\dt=0{,}20$, huit checkpoints jusqu'à $t=16$.
 Les petits $K$ exponentiels conservent du bruit ; la poursuite des
 variantes plus diffusives peut conduire au sur-lissage.}
 \label{fig:curves10}
\end{figure}

Pour \texttt{geometric\_shapes}, $\sigma=0{,}10$, la chaleur atteint
$26{,}7244$~dB à $n=2$, mais seulement $20{,}6655$~dB à $n=80$.
Son SSIM passe de $0{,}8017$ à $n=20$ à $0{,}7138$ à $n=80$.
Le tableau~\ref{tab:time-case} montre que PM peut également sur-lisser :
la rationnelle $K=0{,}20$ passe de $31{,}3253$~dB à $n=5$ à
$21{,}1999$~dB à $n=80$. La stabilité n'implique donc pas que chaque
itération rapproche de la vérité terrain. Le schéma continue à diffuser,
mais les erreurs de structure peuvent prendre le dessus sur le gain
obtenu par réduction du bruit.

\begin{table}[htbp]
 \centering
 \caption{Interaction seuil/temps sur \texttt{geometric\_shapes},
 $\sigma=0{,}10$, aux checkpoints $n=2,5,20,80$ ($t=0{,}4;1;4;16$).
 Moyennes $\pm$ écarts-types échantillonnaux sur dix graines ;
 sous-ensemble explicite des variantes, sans budgets cachés.}
 \label{tab:time-case}
 % Generated from archived summary.csv; rounding applies only to display.
\begin{tabular}{llrll}
\toprule
Méthode & $K$ & $n$ & PSNR (dB) & SSIM \\
\midrule
Chaleur & -- & 2 & $26.7244 \pm 0.0724$ & $0.5828 \pm 0.0054$ \\
Chaleur & -- & 5 & $26.1523 \pm 0.0636$ & $0.7506 \pm 0.0064$ \\
Chaleur & -- & 20 & $23.6482 \pm 0.0432$ & $0.8017 \pm 0.0023$ \\
Chaleur & -- & 80 & $20.6655 \pm 0.0337$ & $0.7138 \pm 0.0009$ \\
PM exp. & 0.2 & 2 & $25.0458 \pm 0.0954$ & $0.3783 \pm 0.0040$ \\
PM exp. & 0.2 & 5 & $31.0611 \pm 0.1795$ & $0.7061 \pm 0.0084$ \\
PM exp. & 0.2 & 20 & $35.7056 \pm 0.3280$ & $0.9378 \pm 0.0026$ \\
PM exp. & 0.2 & 80 & $34.7866 \pm 0.2153$ & $0.9525 \pm 0.0008$ \\
PM rat. & 0.1 & 2 & $23.2927 \pm 0.0720$ & $0.3063 \pm 0.0024$ \\
PM rat. & 0.1 & 5 & $28.6326 \pm 0.1392$ & $0.5643 \pm 0.0067$ \\
PM rat. & 0.1 & 20 & $34.6530 \pm 0.2974$ & $0.9292 \pm 0.0027$ \\
PM rat. & 0.1 & 80 & $28.8937 \pm 0.2417$ & $0.8988 \pm 0.0023$ \\
PM rat. & 0.2 & 2 & $26.9275 \pm 0.0986$ & $0.4623 \pm 0.0041$ \\
PM rat. & 0.2 & 5 & $31.3253 \pm 0.1804$ & $0.7580 \pm 0.0066$ \\
PM rat. & 0.2 & 20 & $26.9975 \pm 0.1527$ & $0.8667 \pm 0.0033$ \\
PM rat. & 0.2 & 80 & $21.1999 \pm 0.0396$ & $0.7315 \pm 0.0010$ \\
\bottomrule
\end{tabular}

\end{table}

La figure~\ref{fig:k-sensitivity} compare les deux conductances aux
mêmes $K$ et aux mêmes temps. La rationnelle $K=0{,}025$, sur le premier
cas à $\sigma=0{,}10$, atteint $23{,}7556$~dB à $n=20$ puis
$36{,}4457$~dB à $n=80$ : un petit seuil peut demander un temps plus long
et n'est pas intrinsèquement défavorable. En revanche, augmenter $K$
favorise une diffusion plus forte et peut nécessiter un arrêt plus tôt.
Une comparaison de seuils détachée de la conductance et du temps serait
donc incomplète.

\begin{figure}[p]
 \centering
 \includegraphics[width=0.98\linewidth]{k_sensitivity_geometric_shapes_sigma0p1.pdf}
 \par\medskip
 \includegraphics[width=0.98\linewidth]{k_sensitivity_ramps_and_edges_sigma0p1.pdf}
 \caption{Sensibilité à $K$ pour $\sigma=0{,}10$ :
 \texttt{geometric\_shapes} en haut, \texttt{ramps\_and\_edges} en bas.
 Deux conductances, $n=5,20,80$ ($t=1,4,16$),
 $K=0{,}025;0{,}05;0{,}10;0{,}20$, $\dx=\dy=1$, $\dt=0{,}20$.
 Points : moyennes sur dix graines ; barres : un écart-type échantillonnal.
 Les références de chaleur sont aux mêmes checkpoints. L'évolution avec
 $K$ change avec le temps : aucun seuil universel n'est identifié.}
 \label{fig:k-sensitivity}
\end{figure}

\subsection{Désaccords entre critères}

Le cas fixé n'est pas isolé. À $\sigma=0{,}025$ sur
\texttt{geometric\_shapes}, une seule itération de chaleur réduit le
PSNR moyen de $32{,}0401$ à $30{,}4691$~dB, mais augmente le SSIM de
$0{,}6646$ à $0{,}8740$. Sur \texttt{ramps\_and\_edges},
$\sigma=0{,}10$, $n=20$, la rationnelle $K=0{,}05$ donne
$32{,}5651\pm0{,}1841$~dB et $0{,}8422\pm0{,}0070$ en SSIM ;
l'exponentielle $K=0{,}20$ donne $32{,}3930\pm0{,}2530$~dB et
$0{,}9209\pm0{,}0023$. Les dispersions invitent à ne pas surinterpréter
le faible écart de PSNR ; les moyennes inversent leur ordre selon le critère,
sans qu'un test de significativité soit effectué ici.

Sur les $42$ couples cas/checkpoint positif, le meilleur réglage selon
le PSNR moyen et celui selon le SSIM moyen diffèrent $13$ fois.
Le compte rendu \texttt{analysis.json} relève aussi $172$ paires
discordantes sur $1\,512$ comparaisons de paires sans égalité.
Ces comptes décrivent la grille, sans test statistique. Enfin,
l'exponentielle $K=0{,}20$ sur \texttt{geometric\_shapes},
$\sigma=0{,}10$, passe entre $n=20$ et $80$ de $35{,}7056$ à
$34{,}7866$~dB, tandis que son SSIM augmente de $0{,}9378$ à $0{,}9525$.
Même le choix d'arrêt peut donc dépendre du critère.

\subsection{Meilleurs réglages observés : un oracle exploratoire}

Le tableau~\ref{tab:oracle-psnr} maximise le PSNR \emph{moyen sur les dix
graines}, parmi les configurations de solveur et checkpoints $n>0$
prédéfinis. C'est un \emph{oracle exploratoire}, puisqu'il utilise
l'image propre. Ce n'est ni un réglage sélectionnable sur une observation
seule ni une performance de généralisation. Il ne maximise pas
séparément chaque graine avant de moyenner : la règle de sélection doit
rester explicite.

\begin{table}[htbp]
 \centering
 \caption{Oracle exploratoire PSNR : maximum de la moyenne sur dix
 graines, parmi la grille figée avec $n>0$. Images et bruits indiqués
 dans chaque ligne ; dispersions échantillonnales. Tout $n=80$ est
 à la frontière temporelle étudiée, pas un optimum global.}
 \label{tab:oracle-psnr}
 % Generated from archived summary.csv; rounding applies only to display.
\begin{tabular}{lllrll}
\toprule
Image / $\sigma$ & Méthode & $K$ & $n$ & PSNR (dB) & SSIM \\
\midrule
Géométrique / 0.025 & PM exp. & 0.05 & 80 & $56.2055 \pm 1.0724$ & $0.9998 \pm 4.10\times10^{-5}$ \\
Géométrique / 0.05 & PM rat. & 0.025 & 80 & $46.1862 \pm 0.6477$ & $0.9967 \pm 0.0004$ \\
Géométrique / 0.1 & PM rat. & 0.05 & 40 & $37.0422 \pm 0.5734$ & $0.9611 \pm 0.0035$ \\
Rampe / 0.025 & PM exp. & 0.05 & 20 & $47.6194 \pm 0.2821$ & $0.9948 \pm 0.0004$ \\
Rampe / 0.05 & PM rat. & 0.025 & 40 & $42.0671 \pm 0.2953$ & $0.9874 \pm 0.0009$ \\
Rampe / 0.1 & PM rat. & 0.025 & 80 & $34.5137 \pm 0.3117$ & $0.9229 \pm 0.0053$ \\
\bottomrule
\end{tabular}

\end{table}

Les choix PSNR/SSIM coïncident dans cinq des six cas. Pour
\texttt{ramps\_and\_edges}, $\sigma=0{,}10$, l'oracle SSIM est en
revanche la rationnelle $K=0{,}05$, $n=40$, alors que l'oracle PSNR
choisit $K=0{,}025$, $n=80$ (tableau~\ref{tab:oracle-difference}).
Plusieurs maxima sont à $n=80$ ; aucun résultat ne permet de les appeler
optima globaux en temps, et aucune expérience supplémentaire n'est
introduite pour prolonger les courbes.

\begin{table}[htbp]
 \centering
 \caption{Deux oracles exploratoires distincts pour
 \texttt{ramps\_and\_edges}, $\sigma=0{,}10$ : sélection selon PSNR ou
 SSIM moyen. Dix graines ; scores et écarts-types issus du CSV,
 $\dx=\dy=1$, $\dt=0{,}20$.}
 \label{tab:oracle-difference}
 % Generated from archived summary.csv; rounding applies only to display.
\begin{tabular}{llrll}
\toprule
Méthode & $K$ & $n$ & PSNR (dB) & SSIM \\
\midrule
PM rat. & 0.025 & 80 & $34.5137 \pm 0.3117$ & $0.9229 \pm 0.0053$ \\
PM rat. & 0.05 & 40 & $34.2340 \pm 0.2965$ & $0.9465 \pm 0.0025$ \\
\bottomrule
\end{tabular}

\end{table}
